\section{Elementary two-class counting and entropy}\label{sec:types}

\subsection{Configurations and ordered profiles}
Give each vertex $a$ distinguished internal stubs and $b$ distinguished
cross stubs. Independently pair the $aN$ internal stubs in each class,
and choose a uniform bijection between the two sets of $bN$ cross stubs.
There are $((aN-1)!!)^2(bN)!$ equally likely configurations. A loop
contributes two to its diagonal adjacency entry and to its neighbor sum.
We use $(-1)!!=1$ throughout.

Let $\mathcal S$ be a fixed finite alphabet. For now its labels belong
to $\Fp$, and local sums are in that field. A plus profile is
\begin{equation}\label{eq:plus-profile}
 \omega=(x;y_1,\ldots,y_a;z_1,\ldots,z_b),\qquad
 \sum_{i=1}^a y_i+\sum_{j=1}^b z_j=\delta x.
\end{equation}
It records the root coordinate, then the coordinates at the partners
of its ordered internal and cross stubs. Define minus profiles in the
same way. There are $S=p^D$ profiles on each side: the root and all
but one internal label determine the last label uniquely.

Let $n^+_\omega,n^-_\omega$ be the profile counts, each summing to $N$,
and write $\pi_\pm=n^\pm/N$. Let $m^\pm_{uv}$ count internal stubs
whose root label is $u$ and requested partner label is $v$. Let
$h^+_{uv}$ count plus cross stubs whose root label is $u$ and requested
minus label is $v$, and let $h^-_{vu}$ count minus cross stubs with
the reverse request. Feasibility requires $m^\pm$ to be symmetric
with even diagonal and $h^+_{uv}=h^-_{vu}$. Write $h$ for the
plus-oriented common table, and put
\begin{equation}\label{eq:tables}
 Q_\pm=\frac{m^\pm}{aN},\qquad R=\frac h{bN},\qquad
 W(m)=\prod_{u<v}m_{uv}!\prod_u(m_{uu}-1)!!.
\end{equation}
The two marginals of $Q_\pm$ are the corresponding root law $\mu_\pm$;
the marginals of $R$ are $\mu_+,\mu_-$.

\begin{lemma}[Exact count]\label{lem:exact-types}
In the configuration model,
\begin{equation}\label{eq:exact-types}
 \E\,|\ker_{\Fp}(A-\delta I)|
 =\sum_{\text{feasible }n^+,n^-}
 \frac{N!}{\prod_\omega n^+_\omega!}
 \frac{N!}{\prod_\omega n^-_\omega!}
 \frac{W(m^+)W(m^-)}{((aN-1)!!)^2}
 \frac{\prod_{u,v}h_{uv}!}{(bN)!}.
\end{equation}
The same formula counts sign eigenvectors if the alphabet is
$\{-1,1\}$ and profiles are restricted to the integer identity
\eqref{eq:plus-profile}.
\end{lemma}
\begin{proof}
The two multinomial factors assign profiles to labeled vertices.
For an assigned internal table, pair the $m_{uv}$ stubs to the
$m_{vu}$ reverse stubs by a bijection when $u<v$, giving $m_{uv}!$.
Pair each diagonal class internally, giving $(m_{uu}-1)!!$.
For the cross table, biject each class to its reverse class, giving
$\prod h_{uv}!$. These choices include loops and multiple edges.
Every configuration and kernel vector determine a unique profile at
each ordered vertex star. Conversely, every compatible choice counted
here satisfies every local eigenvector equation. Thus neither edge
orientation nor root-label multiplicity requires an extra factor.
The reasoning uses only equality of requested endpoint labels and the
local constraint, so it also proves the assertion for integer signs.
\end{proof}

\subsection{Group by aggregate tables, not local types}
The local-type identity is exact but counting all its local types
would introduce a degree-dependent polynomial exponent. We instead
group every ordered profile assignment by its three edge tables.
For a feasible triple $(m^+,m^-,h)$ let $\Omega(m^+,m^-,h)$ be the
set of all ordered arrays of $N$ plus profiles and $N$ minus profiles
having those tables. Thus all assignments of local types to labeled
vertices are already included in $\Omega$. Grouping the same exact
count gives
\begin{equation}\label{eq:exact-fibers}
 \E|\ker_{\Fp}(A-\delta I)|
 =\sum_{m^+,m^-,h}|\Omega(m^+,m^-,h)|
 \frac{W(m^+)W(m^-)}{((aN-1)!!)^2}
 \frac{\prod_{u,v}h_{uv}!}{(bN)!}.
\end{equation}
Empty fibers contribute zero. In particular there is no additional
local-type count or multinomial factor in \eqref{eq:exact-fibers}.

Choose an array uniformly from a nonempty fiber. The fiber is
invariant under independent permutations of vertex positions in each
block and under independent permutations of internal and cross slots
at each vertex. Each such map preserves the local sum and every
aggregate table. Thus all plus vertices have the same slot-symmetric
profile law $\pi_+$, and all minus vertices have the same law $\pi_-$.
Their one-slot edge laws are exactly $Q_+,Q_-,R$ of
\eqref{eq:tables}: invariance makes every position's expected
contribution equal, while the total contribution is fixed by the
fiber. Subadditivity of entropy over the $2N$ vertex profiles gives
\begin{equation}\label{eq:fiber-subadditivity}
 \log|\Omega|\le N\{H(\pi_+)+H(\pi_-)\}.
\end{equation}
This uses no independence between vertices or blocks. The same
statements hold for fibers of integer sign profiles, including fibers
restricted to fixed root counts in each block.

\subsection{Uniform factorial bounds}
All logarithms are natural and $0\log0=0$. Write $H$ for entropy.
The elementary estimates
\begin{equation}\label{eq:factorial}
 r\log r-r\le\log(r!)
 \le r\log r-r+1+\log(r+1),\qquad r\ge0,
\end{equation}
hold with the usual interpretation at zero. For even $r$ use
$(r-1)!!=r!/(2^{r/2}(r/2)!)$.

\begin{lemma}[Zero-safe pairing bounds]\label{lem:factorial-bounds}
Let $q=p(p+1)/2$. For every feasible internal table of total $T=aN$,
and every cross table of total $bN$,
\begin{align}
 \frac{W(m)}{(T-1)!!}
 &\le e^{q+1}(T+1)^{q+1}
       \exp\!\left[-\frac T2H(m/T)\right],\label{eq:pair-bound}\\
 \frac{\prod h_{uv}!}{(bN)!}
 &\le e^{p^2}(bN+1)^{p^2}\exp[-bNH(R)].\label{eq:cross-bound}
\end{align}
The multinomial factor for a profile law $\pi$ is at most $e^{NH(\pi)}$.
\end{lemma}
\begin{proof}
For an off-diagonal class the leading term in $\log(m_{uv}!)$ is
$m_{uv}\log m_{uv}-m_{uv}$. For a diagonal class the leading term
in $\log((m_{uu}-1)!!)$ is
$(m_{uu}/2)\log m_{uu}-m_{uu}/2$. Summing, their leading terms are
$\tfrac12\sum_{u,v}m_{uv}\log m_{uv}-T/2$.
The denominator has leading term $(T/2)\log T-T/2$.
There are $q$ numerator classes and one denominator. Applying
\eqref{eq:factorial}, or its double-factorial consequence, bounds
the upper remainder by $(q+1)(1+\log(T+1))$. This proves
\eqref{eq:pair-bound}, including empty classes.
For the cross ratio the leading terms leave $-bNH(R)$ and the
$p^2$ numerator factorials cost at most
$p^2(1+\log(bN+1))$; a lower factorial bound suffices in the denominator.
Finally the probability of a type under its own product law is at
most one, which gives the multinomial entropy bound.
\end{proof}

Consequently each summand of \eqref{eq:exact-types} is bounded by a
fixed polynomial factor times $e^{N\Phi}$, where
\begin{equation}\label{eq:Phi}
 \Phi=H(\pi_+)+H(\pi_-)
       -\frac a2\{H(Q_+)+H(Q_-)\}-bH(R).
\end{equation}
All these bounds are uniform at the boundary of type space. No central
limit theorem, nondegenerate Hessian, or asymptotic saddle constant
will be needed.

\subsection{The entropy cancellation}
Independently average each profile law over permutations of its $a$
internal positions and its $b$ cross positions. The constraint,
aggregated tables, and root laws are preserved. Profile entropy can
only increase. It suffices to upper-bound the resulting symmetrized
$\Phi$.

For the plus law write its random profile as
$(X;Y_1,\ldots,Y_a;Z_1,\ldots,Z_b)$. Internal table symmetry and cross
table consistency imply
\[
 \mathcal L(X)=\mathcal L(Y_i)=\mu_+,
 \qquad \mathcal L(Z_j)=\mu_-.
\]
The analogous minus profile has root law $\mu_-$, internal law
$\mu_-$, and cross law $\mu_+$. The constraint determines $Y_a$ from
the root and the retained neighbors. Thus
\begin{equation}\label{eq:subadd}
 H(\pi_+)=H(\mu_+)+(a-1)H(Y_1\mid X)+bH(Z_1\mid X)-J_+,
 \qquad J_+\ge0,
\end{equation}
where $J_+$ is conditional total correlation among the retained slots.
Define $J_-$ in the same way. Put
\[
 I_+=2H(\mu_+)-H(Q_+),\quad
 I_-=2H(\mu_-)-H(Q_-),\quad
 I_\times=H(\mu_+)+H(\mu_-)-H(R).
\]
These are the internal and cross mutual informations and are nonnegative.
Substitution in \eqref{eq:Phi} gives the exact identity
\begin{equation}\label{eq:entropy-identity}
 -\Phi=\left(\frac a2-1\right)(I_++I_-)
          +bI_\times+J_++J_-.
\end{equation}
In particular $\Phi\le0$ whenever $a\ge2$, over every finite field.
For an unsymmetrized law its original exponent is at most the
symmetrized exponent, so this upper bound remains valid.

Apply this entropy inequality to the uniform-fiber laws in
\eqref{eq:fiber-subadditivity}. We obtain
\begin{equation}\label{eq:fiber-bound}
 |\Omega|\le
 \exp\!\left[N\left\{\frac a2(H(Q_+)+H(Q_-))+bH(R)\right\}\right].
\end{equation}
Multiplying by the pairing bounds cancels this exponential exactly.
There are at most $(aN+1)^{2p^2}(bN+1)^{p^2}$ table triples, bounding
each entry independently and ignoring symmetry, parity, and totals.
Equation~\eqref{eq:exact-fibers} therefore gives
\begin{equation}\label{eq:explicit-moment}
 \E\,|\ker_{\Fp}(A-\delta I)|
 \le e^{p(p+1)+2+p^2}
 (aN+1)^{p(p+1)+2+2p^2}(bN+1)^{2p^2}.
\end{equation}
Its total power of $N$ is $5p^2+p+2$, independent of the fixed degrees.
In particular,
\begin{equation}\label{eq:explicit-binary-moment}
 \E\,2^{\dim_{\F_2}\ker(A-\delta I)}
 \le e^{12}(aN+1)^{16}(bN+1)^8
 =O_{a,b}(N^{24}).
\end{equation}
The fiber step avoids any prefactor for enumerating local profiles;
its cardinality was bounded directly by entropy, not by a sum of
type-by-type estimates.

\subsection{Conditioning gives exactly the uniform-simple model}
Every simple labeled internal $a$-regular graph has $(a!)^N$
distinguished-stub realizations. Every simple labeled $b$-regular
bipartite graph has $(b!)^{2N}$ cross-bijection realizations.
Conditioning the three independent factors on simplicity therefore
gives exactly the three independent uniform-simple laws in our model.
At fixed degrees their joint simplicity probability is bounded below
by a positive constant along all admissible sizes. This follows, for
example, from the ordinary and bipartite configuration criteria in
Janson~\cite[Theorems 1.1 and 5.1]{JansonSimple}; the latter theorem
credits the bipartite criterion to Blanchet and Stauffer.
For a nonnegative random variable $Z$ and the joint simplicity event
$\mathcal S$, $\E[Z\mid\mathcal S]\le\E Z/\Prob(\mathcal S)$.
Applying this to \eqref{eq:explicit-moment} proves the same elementary
polynomial estimate in the simple law. The stronger bounded moment in
Theorem~\ref{thm:general-moment} follows directly from
Theorem~\ref{thm:frame-bounded-moment}, with two equal classes, frame
$Q=\left(\begin{smallmatrix}a&b\\b&a\end{smallmatrix}\right)$,
and scalar target $\delta$. The elementary count above supplies a
separate check and the explicit prefactor used in the Boolean proof.
Both transfers use actual positive-probability simplicity conditioning,
rather than qualitative contiguity of different graph laws.
